\subsection{Lokale Beleuchtungsmodelle}
\label{subsec:local_illumination_models}

Lokale Beleuchtungsmodelle aggregieren Daten von benachbarten, eben lokalen,
Oberflächen. Diese Modelle sind in deren Umfang allerdings limitiert, da sie
normalerweise nur Lichtquellen sowie die Orientierung einer Oberfläche
einbeziehen. Sie ignorieren dabei aber die globale Umgebung, in welcher sich
eine Oberfläche befindet.  Dies ist dadurch bedingt, dass die traditionell
verwendeten Algorithmen zur Berechnung der Sichtbarkeit von Oberflächen, über
keine globalen Daten verfügen. Das bekannteste lokale Beleuchtungsmodell ist
das \textbf{Phong}-Beleuchtungsmodell, welches 1975 von Phong Bui-Tong
entwickelt wurde. Es beschreibt nicht-perfekte Reflektoren wie zum Beispiel einen
Apfel~\cite{foley_computer_1996}[Kapitel 16, Seite 729].

\begin{figure}[H]
    \centering
    \includegraphics[width=0.6\textwidth]{img/phong_illumination_model.pdf}
    \caption{Illustration des Phong-Beleuchtungsmodelles\protect\footnotemark}\label{fig:phong_illustration}
\end{figure}
\footnotetext{Eigene Darstellung mittels Geogebra, angelehnt an~\cite{foley_computer_1996}[Kapitel 16, Seite 731, Abbildung 16.12]}

Es beschreibt die reflektierte (Licht-) Intensität als Zusammensetzung aus der ambienten, der diffusen und der ideal spiegelnden Reflexion einer Oberfläche:

\begin{gather}
    I = I_{\text{ambient}} + I_{\text{diffuse}} + I_{\text{specular}} + I_{\text{emissive}}
\end{gather}

oder mathematisch ausgedrückt:

\begin{gather}
    I(\vv{V}) = k_{a} \cdot L_{a} +
                k_{d} \displaystyle\sum_{i=0}^{n - 1} L_{i} \cdot (\vv{S_{i}} \cdot \vv{N}) +
                k_{s} \displaystyle\sum_{i=0}^{n - 1} L_{i} \cdot {(\vv{R_{i}} \cdot \vv{V})}^{k_{e}}
\end{gather}

wobei gilt:

\begin{itemize}
    \item $I(\vv{V})$:              Die reflektierte (Licht-) Intensität in Richtung des Vektors $\vv{V}$
    \item $n$:                      Anzahl Lichtquellen
    \item $k_{a} \cdot L_{a}$:      Ambiente Komponente des
                                    Beleuchtungsmodelles. Mittels diesem Faktor
                                    wird versucht allem indirekten Licht der
                                    Szene gerecht zu werden. Bei $k_{a}$
                                    handelt es sich um eine Konstante, welche
                                    den ambienten Lichtanteil $L_{a}$ skaliert
    \item $k_{d}$:                  Konstante für die diffuse Komponente des
                                    reflektierten Lichtes, basierend auf der
                                    Wellenlänge bzw. Frequenz
    \item $\vv{S_{i}}$:             Richtung, in welcher das Licht der $i$-ten
                                    Lichtquelle ankommt, normalisierter
                                    Einheitsvektor
    \item $\vv{N}$:                 Einheitsnormale der Oberfläche
    \item $k_{s}$:                  Koeffizient der spiegelenden Komponente,
                                    basierend auf der Wellenlänge bzw. Frequenz
    \item $\vv{R_{i}}$:             Richtung, in welcher das Licht der $i$-ten
                                    Lichtquelle reflektiert wird,
                                    normalisierter Einheitsvektor
    \item $\vv{V}$:                 Blickrichtung des Betrachters bzw.\ der
                                    Kamera
    \item $k_{e}$:                  Exponent, welcher von der Rauheit bzw.
                                    Reflektion der Oberfläche abhängt
    \item $L_{i}$                   Licht- bzw.\ Farbintnsität der $i$-ten
                                    Lichtquelle
\end{itemize}

Meist wird der emissive Term bewusst weggelassen, da dieser meist eher für
Spezialeffekte statt für die Beleuchtung ``normaler'' Objekte benutzt wird.

Der reflektive Vektor $R_{i}$ ist gegeben durch

\begin{gather}
    R_{i} = S_{i} + 2(S_{i} \cdot N)N
\end{gather}

Damit die Energieerhaltung gewährleistet ist, muss weiter $k_{d} + k_{s} < 1$
gelten. Der Winkel zwischen $\vv{R}$ und $\vv{V}$ wird mittels $\cos{\alpha}$
ermittelt.
